\chapter{类型、表达式与控制流：写出第一个行情筛选器}

\begin{introduction}
  \item 把 Python 的“名称绑定对象”心智模型切换为 C++ 的“声明对象并初始化”
  \item 选择基本类型、定宽整数、字面量和 \texttt{auto}，并识别窄化
  \item 用算术、比较、布尔表达式和显式转换计算行情指标
  \item 用 \texttt{if}/\texttt{else}、\texttt{switch}、\texttt{for} 和 \texttt{while} 控制执行路径
  \item 从标准输入读取固定行情，筛选有效买单并汇总名义金额
\end{introduction}

\chaptermeta{已完成第 1 章，能在 Linux/WSL 中编译、链接并运行单文件程序；本章不要求会自定义函数、容器、类、引用或指针。}{Python 变量是可重新绑定的名称；C++ 声明会创建具有固定类型和作用域的对象，复制一个基础类型对象会得到独立的值。}{先区分编译诊断与运行结果，再检查初始化、整数除法、隐式转换、分支边界和循环是否终止。}

\section{本章任务：从五行行情得到一个确定答案}

研究原型里，一条筛选可能只是几行 DataFrame 表达式；面试白板题或低延迟数据入口却常要求你明确说明：每个字段是什么类型，坏数据走哪条分支，循环何时停止，累计值如何更新。本章的输出任务只使用基础语法，不借助尚未教学的自定义函数、结构体或容器。

程序从标准输入读取一个行数和五行数值行情。每行依次是“证券编号、方向代码、价格、数量”；方向 \texttt{1} 表示买，\texttt{2} 表示卖，其他代码无效。程序只选中编号、价格、数量均为正的买单，然后打印选中数、拒绝数、名义金额总和与每条选中记录的平均名义金额。

完成本章后，下面的输入必须得到逐字符一致的核心数值结果：

\begin{lstlisting}[numbers=none]
5
101 1 100.5 10
102 2 50.0 4
103 1 25.0 0
104 1 10.25 3
105 9 12.0 5
\end{lstlisting}

\begin{lstlisting}[numbers=none]
selected=2 rejected=3 buy_notional=1035.75 average=517.875
\end{lstlisting}

本章以可执行能力验收：读者应在不复制最终答案的前提下重写程序、解释每个对象的状态变化，并用给定数据复现输出；仅仅“看懂语法”还不够。

\section{对象、类型与初始化}

\subsection{场景：一行数据需要四个有边界的对象}

读入一行行情时，程序需要保存证券编号、方向、价格和数量。如果这些值没有确定的初始状态，后面的筛选就可能读取垃圾值；如果价格被随意塞进整数，小数部分会悄悄丢失。C++ 因而要求先回答“创建什么类型的对象”，再讨论计算。

\subsection{语法规则与最小实验}

最常见的声明由“类型、对象名、初始化器、分号”组成。下面给出完整程序，会参与 C++20 构建并由 CTest 检查输出：

\lstinputlisting[caption={对象、字面量与独立复制}]{code/ch02/objects_and_literals.cpp}

\texttt{int} 保存整数，\texttt{double} 保存双精度浮点数，\texttt{bool} 只有 \texttt{true} 和 \texttt{false} 两种逻辑状态，\texttt{char} 保存一个字符。\texttt{std::int64\_t} 来自 \texttt{<cstdint>}，在提供它的平台上恰有 64 位；当数量或时间戳需要明确位宽时，它比“希望 \texttt{long} 足够大”更可检查。

花括号 \texttt{\{\}} 同时标记对象从这里开始生命周期，并拒绝一部分会丢信息的窄化。空花括号执行值初始化：\texttt{int selected\{\};} 得到零，\texttt{bool is\_buy\{\};} 得到 \texttt{false}。相反，局部基础类型 \texttt{int selected;} 没有可靠的业务初值；读取它属于程序错误，不能把所得值视为“随机但可用”。

字面量本身也有类型。\texttt{1} 通常是 \texttt{int}，\texttt{1.0} 是 \texttt{double}，\texttt{true} 是 \texttt{bool}，单引号中的 \texttt{'X'} 是字符字面量，双引号中的 \texttt{"side="} 是字符串字面量。\texttt{'\textbackslash n'} 表示换行字符，它与包含两个可见字符的 \texttt{"\textbackslash\textbackslash n"} 不同。因此 \texttt{double price\{100\};} 可以精确表示该值，而 \texttt{int price\{100.5\};} 会在编译阶段被拒绝。

为了让本章不依赖读者猜测程序外壳，这里逐项说明：\texttt{\#include <cstdint>} 与 \texttt{\#include <iostream>} 让当前源文件看见所需标准库声明；\texttt{int main()} 定义程序入口，\texttt{int} 是返回状态的类型，空圆括号表示这个版本不接收命令行参数，随后的一对花括号包住入口函数体。\texttt{std::cout} 是标准输出流，\texttt{<<} 依次把右侧的字符串字面量和对象值送入流。\texttt{std::} 指明标准库名称所在的命名空间。最后一个右花括号结束 \texttt{main}；走到这里等价于成功返回状态 0。第 3 章才会系统教学如何定义和调用自己的函数。

\subsection{逐步观察：Python 名称不等于 C++ 对象}

在 Python 中，赋值把名称绑定到对象；两个名称可以指向同一个可变列表：

\begin{lstlisting}[language=Python]
a = [10]
b = a
a[0] = 11
print(b[0])          # 11
\end{lstlisting}

在上面的完整 C++ 程序中，\texttt{original} 与 \texttt{copied} 是两个独立的 \texttt{int} 对象。\texttt{copied\{original\}} 用 \texttt{original} 当时的值初始化新对象，之后把 \texttt{original} 赋值为 11，不会改动 \texttt{copied}；可观察输出同时包含 \texttt{changed=11} 与 \texttt{copied=10}。

等号在对象创建之后出现时是赋值，不是初始化。对象的类型仍是 \texttt{int}，不能像 Python 名称那样下一行重新绑定成字符串。共享和借用当然可以在 C++ 中表达，但要用后续章节才教学的引用、指针或拥有型对象，不能从普通按值声明中猜出来。

\begin{pythonbridge}
两种语言都允许把已有值用于新变量；边界在于 Python 名称通常保存对象引用，而本章的 C++ 基础类型声明创建独立对象。Python 类型标注通常不改变运行期绑定规则，C++ 对象类型则直接限制初始化、运算和存储表示。
\end{pythonbridge}

\subsection{失败诊断与立即练习}

故意失败文件 \path{code/ch02/failures/narrowing.cpp} 包含如下核心语句：

\begin{lstlisting}
int rounded_price{100.5};  // intentional compile failure
\end{lstlisting}

用 \texttt{g++ -std=c++20 code/ch02/failures/narrowing.cpp} 单独触发它。预期诊断类别是“列表初始化拒绝窄化转换”，具体英文措辞可随 GCC 版本变化。根因是从 \texttt{double} 字面量到 \texttt{int} 会丢掉小数；正确修复需要先决定业务舍入规则，再显式转换，换成圆括号只会绕开检查。

立即练习：分别编译 \texttt{double price\{10\};} 与 \texttt{int quantity\{10.0\};}。前者的整数常量 10 能被 \texttt{double} 精确表示，可以通过；后者从浮点类型转为整数，即使写成 10.0 也属于列表初始化禁止的窄化，必须被拒绝。再把它改成 10.5，说明同一诊断规则，而不是误记成“没有小数部分就可以”。

\section{表达式、\texttt{auto} 与显式转换}

\subsection{场景：价格乘数量还不够}

名义金额是价格与数量的乘积，平均值还要除以选中记录数。这里会同时遇到运算符优先级、整数除法、结果类型和除零边界。Python 常把数值提升留给运行时；C++ 会先根据操作数类型决定表达式类型。

\subsection{语法规则与可运行路径}

算术运算符包括 \texttt{+}、\texttt{-}、\texttt{*}、\texttt{/} 和整数余数 \texttt{\%}。乘除先于加减；括号既能改变顺序，也能让审阅者看清意图。比较运算 \texttt{<}、\texttt{<=}、\texttt{>}、\texttt{>=}、\texttt{==}、\texttt{!=} 产生 \texttt{bool}。

下面的完整程序由构建目标 \texttt{ch02\_expressions} 编译，并输出三个可独立核算的结果：

\lstinputlisting[caption={算术、类型推导与显式转换}]{code/ch02/expressions.cpp}

\texttt{auto} 让编译器从初始化表达式推导对象类型；它不是 Python 式动态类型。这里推导结果仍固定为 \texttt{double}。必须同时写初始化器，不能写孤立的 \texttt{auto notional;}。当类型本身承载业务信息时仍应显式写出，例如数量使用 \texttt{std::int64\_t}，不要为了少打字全部改成 \texttt{auto}。

\texttt{static\_cast<double>(quantity)} 明确表达“此处要以浮点数参与计算”。对本例的乘法，即使不写转换，\texttt{double} 价格也会促使数量转换为 \texttt{double}；保留显式转换是为了把边界写给读者，并为随后的除法建立习惯。程序依次得到 \texttt{notional=1005}、\texttt{average=502.5} 和 \texttt{integer\_division=2}；最后一个结果来自两个 \texttt{int} 的 \texttt{5 / 2}。

复合赋值 \texttt{buy\_notional += notional;} 等价于按该运算规则更新左侧对象；\texttt{++selected;} 把整数加一。二者都会修改左侧对象；如何声明只读对象留到第 3 章。

\begin{pythonbridge}
Python 的 \texttt{/} 总是产生浮点除法，\texttt{//} 才是向下取整除法；C++ 的 \texttt{/} 取决于两侧类型，两个整数相除先得到整数结果。\texttt{auto} 只省略静态类型拼写，不允许对象之后换类型。
\end{pythonbridge}

\subsection{失败诊断与立即练习}

若平均值意外变成 \texttt{602}，先打印或说明除法两侧的类型，不要先怀疑格式化。另一个常见失败是选中数为零时仍做除法；浮点除零可能得到无穷而继续污染结果，整数除零则是更严重的运行期错误。正确做法是在控制流中显式保护分母。

立即练习：分别预测并运行 \texttt{7 / 3}、\texttt{7.0 / 3}、\texttt{static\_cast<double>(7) / 3}。写出实际结果，并指出哪一步发生类型转换。

\section{条件表达式、\texttt{if} 与 \texttt{switch}}

\subsection{场景：有效买单必须同时通过多道门}

本章规则要求证券编号、价格、数量都为正，而且方向必须是买。任何一项失败都拒绝该行。把条件拆清楚比写一个“聪明”的长表达式更适合面试解释和生产审阅。

\subsection{布尔规则与短路求值}

逻辑与写作 \lstinline|&&|，逻辑或写作 \lstinline!||!，逻辑非写作 \texttt{!}。\lstinline|a && b| 只有两侧都为真才为真，并且当 \texttt{a} 已为假时不会求值 \texttt{b}；\lstinline!a || b! 在 \texttt{a} 已为真时不会求值 \texttt{b}。这种短路行为既能表达保护条件，也意味着右侧可能根本不执行，不能偷偷依赖它产生副作用。

完整分支实验同时包含布尔表达式、\texttt{switch} 和 \texttt{if}/\texttt{else}，给定内置行情时输出 \texttt{selected}：

\lstinputlisting[caption={离散方向分派与有效性筛选}]{code/ch02/branches.cpp}

\texttt{if} 后的圆括号放条件；条件为真执行第一对花括号，为假执行 \texttt{else} 的花括号。即使分支只有一行，本书也保留花括号，避免后来新增语句时出现视觉缩进与真实控制流不一致。

方向代码是少量离散整数，适合用清单中的 \texttt{switch} 明确列出分支。\texttt{case} 标签必须是可在编译期确定的离散值。\texttt{break} 离开当前 \texttt{switch}；漏写时会继续执行下一标签，这叫贯穿，除非有明确意图和注释，否则通常是缺陷。\texttt{default} 接住未列出的代码，使新出现的脏数据不会悄悄被当成买单。

\begin{pythonbridge}
Python 的 \texttt{if}/\texttt{elif}/\texttt{else} 与 C++ 分支目标相似，但 C++ 用圆括号和花括号界定结构，并要求条件可转换为布尔值。C++ \texttt{switch} 不是 Python \texttt{match} 的完整对应物：这里只把它当作离散整数分派，不期待结构模式匹配。
\end{pythonbridge}

\subsection{失败诊断与立即练习}

\texttt{if (side\_code = 1)} 是赋值，不是比较；对象被改成 \texttt{1}，条件随后为真。高警告级别通常会提示这一可疑写法。修复为 \texttt{==}，并保持 \texttt{-Wall -Wextra -Wpedantic}。另一个边界是把 \texttt{price >= 0.0} 写成有效规则，这会错误接纳零价格；先从业务条件写出真值表，再翻译成表达式。

立即练习：把有效规则改成“买单且价格至少 10，数量在 1 到 100（含端点）”。用恰好位于 10、1、100 的三条记录验证端点，再用价格 9.99 验证拒绝路径。

\section[循环与作用域]{\texttt{for}、\texttt{while} 与作用域}

\subsection{场景：已知行数与未知次数是两种循环}

输入开头已经给出行数，适合用 \texttt{for} 精确处理五次；风险预算能接纳多少笔交易要到条件失败时才知道，适合用 \texttt{while}。选择循环形式是在表达终止依据，不只是个人风格。

\subsection{计数循环与对象可见范围}

计数循环也有独立、可构建的最小程序。它检查 0 到 4 的余数，选中三个偶数行：

\lstinputlisting[caption={计数循环与块作用域}]{code/ch02/for_scope.cpp}

圆括号内由两个分号分隔的三个部分依次是：只执行一次的初始化 \texttt{int row\{\}}；每轮开始前检查的条件 \texttt{row < 5}；每轮末尾执行的更新 \texttt{++row}。\texttt{row} 因而依次为 0、1、2、3、4，共五轮；若误写成 \texttt{row <= 5} 就会多处理一轮。

每轮先用 \texttt{row \% 2} 计算余数，并创建本轮的 \texttt{remainder}。余数为零时，\texttt{if} 分支把循环外的 \texttt{selected} 加一，因此第 0、2、4 轮共更新三次并输出 \texttt{selected=3}。花括号还建立作用域：\texttt{remainder} 每轮重新创建并在该轮末尾销毁，\texttt{selected} 则跨轮保存累计值；循环结束后也不能访问初始化部分声明的 \texttt{row}。把对象放在“足够用的最小作用域”能减少误用，也让读者更容易追踪状态。

\subsection{条件循环的完整练习答案}

当循环次数由状态决定时，把条件放在 \texttt{while} 后。下面的完整答案从 3200 元风险预算中反复扣除每笔 750 元，只有预算足够时才进入下一轮：

\lstinputlisting[caption={用 while 计算可接纳交易数}]{code/ch02/exercises/risk_budget_solution.cpp}

初始预算 3200；四轮后剩 200，第五轮开始前 \texttt{200 >= 750} 为假，因此接纳数为 4、剩余预算为 200。循环体必须让条件朝终止方向变化；若忘记扣减预算，就会形成无限循环。

\begin{pythonbridge}
两种语言的 \texttt{while} 都在每轮前检查条件。Python 常写 \texttt{for row in range(n)}；C++ 的经典 \texttt{for} 把初始化、条件和更新显式放在一处。本章还没有容器，所以暂不使用范围 \texttt{for}，它会在第 4 章教学。
\end{pythonbridge}

\subsection{失败诊断与立即练习}

循环结果差一条时，手工列出“进入本轮前的 \texttt{row}、条件结果、本轮后的 \texttt{row}”三列，比盯着代码更快。程序长时间不结束时，检查每条分支是否都会更新终止状态，以及输入失败后是否还在重复使用旧值。

立即练习：把风险预算改为 2999、每笔风险改为 1000。先预测循环次数与余数，再运行验证，结果应为 \texttt{accepted\_trades=2 remaining\_budget=999}。

\section{整合：读取、筛选并汇总固定行情}

下面是本章输出任务的完整源码。它只组合已经教学的对象、初始化、表达式、转换、分支、循环和作用域。这里补齐输入与错误流规则，不要求读者从旧版第 1 章猜语法：\texttt{std::cin >> row\_count} 把以空白分隔的文本转换后写入对象，多个 \texttt{>>} 可以从左到右连续提取；输入失败时流状态在条件中为假，前缀 \texttt{!} 将它取反。\texttt{std::cerr} 用于错误消息，正常结果仍送到已经教学的 \texttt{std::cout}。

\lstinputlisting[caption={固定行情筛选与汇总}]{code/ch02/market_filter.cpp}

\subsection{按执行顺序走一遍}

\begin{enumerate}
  \item \texttt{row\_count\{\}} 先得到零；输入成功后变成 5。若读取失败或行数为负，程序向错误流报告并返回非零状态。
  \item 三个累计对象从零开始。它们声明在循环外，因此每轮的更新会保留到下一轮。
  \item 每轮创建四个字段对象并读取一行。第 1 行方向为买、各字段为正，行名义金额为 $100.5\times10=1005$。
  \item 第 2 行是卖，第 3 行数量为零，第 5 行方向代码未知，三者都走拒绝分支。第 4 行贡献 $10.25\times3=30.75$。
  \item 循环结束时总额为 $1005+30.75=1035.75$，选中数为 2。分母保护通过后，平均值为 $1035.75/2=517.875$。
  \item 输出流依次写出标签与对象值，所以得到本章开头约定的确定字符串。
\end{enumerate}

在 Linux/WSL 中，从项目根目录运行：

\begin{lstlisting}[language=bash,numbers=none]
cmake -S . -B build
cmake --build build --target ch02_market_filter
./build/ch02_market_filter < code/ch02/data/market_rows.txt
ctest --test-dir build -R ch02_ --output-on-failure
\end{lstlisting}

Windows 的生成器可能把程序放到配置子目录，或在文件名后加 \texttt{.exe}；源码行为和测试期望不变。CTest 的预期字符串是独立写入的已知结果，不调用被测程序的算法重新计算自己。

\section{自测、练习与面试表达}

先不看答案，逐项给出定义、使用时机、代价或失败例子：

\begin{enumerate}
  \item 初始化和赋值分别发生在对象生命周期的什么位置？\texttt{int x;} 与 \texttt{int x\{\};} 对局部对象有何差异？
  \item 为什么 \texttt{auto} 不是动态类型？何时宁可显式写 \texttt{std::int64\_t}？
  \item 解释 \texttt{5 / 2} 与 \texttt{5.0 / 2} 的结果，并给出一个平均值被整数截断的失败例子。
  \item \lstinline|&&| 的短路规则是什么？它怎样保护右侧表达式，又可能怎样隐藏副作用？
  \item \texttt{switch} 中漏写 \texttt{break} 会怎样？为什么仍应保留 \texttt{default}？
  \item 分别给出适合 \texttt{for} 和 \texttt{while} 的终止依据，并说明一个无限循环的根因。
\end{enumerate}

编码练习：

\begin{enumerate}
  \item 从空文件开始重写行情筛选器，不复制清单。先只输出 \texttt{selected} 和 \texttt{rejected}，再加入总额与平均值；使用给定输入达到本章约定输出。完整可运行答案是 \path{code/ch02/market_filter.cpp}。
  \item 写一个 \texttt{while} 程序：风险预算足够时接纳一笔固定风险交易，最后打印接纳数和余数。用 3200 与 750 得到 4 和 200。完整可运行答案是 \path{code/ch02/exercises/risk_budget_solution.cpp}。
  \item 单独编译窄化失败实验，记录失败阶段和诊断类别；然后提出两种修复：改变目标类型，或在明确舍入政策后显式转换。不要把关闭警告当成修复。
\end{enumerate}

\begin{interviewcheck}
用 90 秒回答：“Python 名称绑定与 C++ 对象初始化有何不同？为什么花括号、显式类型和短路条件对行情入口有价值？”答案至少包含一个可观察输出、一个编译期失败和一个循环边界，而不是只罗列术语。
\end{interviewcheck}

\section{本章小结}

C++ 程序从具有固定类型和作用域的对象开始；初始化建立初始状态，表达式根据操作数类型产生结果，显式转换标出有意的类型边界。\texttt{if} 和 \texttt{switch} 选择路径，\texttt{for} 和 \texttt{while} 重复路径，花括号同时限定控制结构与对象可见范围。你现在已经能只用基础语法完成一个有输入、有脏数据分支、有循环汇总且可测试的量化小程序。后续章节会在这一基础上加入函数、资源生命周期和标准容器。
